
\documentclass[12pt,aspectratio=169]{beamer}
\input{header_beamer}
\usetheme{Boadilla}
\usecolortheme{beaver}
\setbeamertemplate{page number in head/foot}{}
\setbeamertemplate{navigation symbols}{}

\title{Lecture-06: Tranformation of Random Vectors}
\date{}

\begin{document}
\pagestyle{empty}
\maketitle
\begin{frame}{Functions of random variables}
\begin{Definition}
\begin{itemize}
\item<1-> Borel measurable sets on a space $\R^n$ is denoted by $\sB(\R^n)$ and generated by the collection $(\pi_i^{-1}(-\infty,x]: x \in \R, i \in [n])$. 
\item<2-> A function $g: \R^n\to \R^m$ is called \textbf{Borel measurable} function, 
if $g^{-1}(B_m) \in \sB(\R^n)$ for any $B_m \in \sB(\R^m)$. 
\end{itemize}
\end{Definition}
\begin{prop}<3->
Consider a random variable $X: \Omega \to \R$ defined on the probability space $(\Omega, \sF, P)$. 
Suppose $g: \R \to \R$ is function such that $g^{-1}(-\infty, x] \in \sB(\R)$, 
then $g(X)$ is a random variable. 
\end{prop}
\end{frame}

\begin{frame}{Functions of random variables}
\begin{Proof}
\begin{itemize}
    \item<1-> Represent $g(X)$ by a map $Y: \Omega \to \R$ such that $Y(\omega) \triangleq (g \circ X)(\omega)$ for all outcomes $\omega \in \Omega$. 
    \item<2-> Check that for any half open set $B_x = (-\infty,x]$, we have $Y^{-1}(B_x) = (X^{-1}\circ g^{-1})(B_x)$. 
    \item<3-> Since $g^{-1}(B_x) \in \sB(\R)$, it follows that $Y^{-1}(B_x) \in \sF$ by the definition of random variables.
\end{itemize}
\end{Proof}
\end{frame}

\begin{frame}{Functions of random variables}
\begin{Example}[Monotonically increasing function of random variables] 
\begin{itemize}
\item <1-> Let $g:\R \to \R$ be a monotonically increasing function, 
then $g^{-1}(-\infty, x] %= (g^{-1}(-\infty), g^{-1}(x)] 
\in \sB(\R)$ for all $x \in \R$. 
\item <2-> Consider a random variable $X: \Omega \to \R$ defined on the probability space $(\Omega, \sF, P)$, 
then $Y \triangleq g(X)$ is a random variable with distribution function
\EQ{
F_Y(y) = P\set{g(X) \le y} = P\set{X \le g^{-1}(y)} = F_{X}(g^{-1}(y)). 
}
\item <3-> Here, $g^{-1}(y)$ is the functional inverse, and not inverse image as we have been seeing typically. 
We can think $g^{-1}(y) = g^{-1}\set{y}$, though this inverse image has at most a single element since $g$ is monotonically increasing. 
\end{itemize}
\end{Example}
\end{frame}

\begin{frame}{Functions of random variables}
\begin{Example}[Monotonically decreasing function of random variables]
\begin{itemize}
\item <1-> Consider a positive random variable $X: \Omega \to \R_+$ defined on a probability space $(\Omega,\sF,P)$.    
\item <2-> Let $g: \R_+ \to \R_+$ be such that $g(x) = e^{-\theta x}$ for all $x \in \R_+$ and some $\theta > 0$.  
\item <3-> Then, $g$ is monotonically decreasing in $X$ and $x = g^{-1}(y) = -\frac{1}{\theta}\ln y$. 
\item <4-> This implies that $g^{-1}(-\infty, y] = [-\frac{1}{\theta}\ln y, \infty) \in \sB(\R_+)$ for all $y \in \R_+$. 
\item <5-> Thus $g$ is a measurable function, and $Y = g(X)$ is a random variable.
\end{itemize}
\end{Example}
\end{frame}

\begin{frame}{Functions of random variables}
\begin{prop}[Independence of function of random variables] 
\begin{itemize}
\item <1-> Let $g:\R \to \R$ and $h: \R \to \R$ be functions such that $g^{-1}(-\infty, x]$ and $h^{-1}(-\infty, x]$ are Borel sets for all $x \in \R$. \\
\item <2-> Consider independent random variables $X$ and $Y$ defined on the probability space $(\Omega, \sF, P)$, 
then $g(X)$ and $h(Y)$ are independent random variables. 
\end{itemize}
\end{prop}
\end{frame}

\begin{frame}{Functions of random variables}
\begin{Proof}
\begin{itemize}
\item <1-> For any $u,v \in \R$, we can define inverse images $A_g(u) \triangleq g^{-1}(-\infty,u]$ and $A_h(v) \triangleq h^{-1}(\infty, v]$. 
\item <2-> Since $g,h$ are Borel measurable, we have $A_g(u), A_h(v) \in \sB(\R)$.   
\item <3-> We can write the following outcome set equality for the joint event 
\EQ{
\set{g(X) \le u}\cap\set{h(Y) \le v} = \set{X \in g^{-1}(-\infty, u]}\cap\set{Y \in h^{-1}(-\infty,v]}}
\EQ{= X^{-1}(A_g(u))\cap Y^{-1}(A_h(v)) \in \sF.
}
\item <4-> Since $X$ and $Y$ are independent random variables, it follows that $X^{-1}(A_g(u))$ and $Y^{-1}(A_h(v))$ are independent events, 
and the result follows. 
\end{itemize}
\end{Proof}
\end{frame}

\begin{frame}{Function of random vectors}
\begin{prop}
\begin{itemize}
\item <1-> Consider a random vector $X: \Omega \to \R^n$ defined on the probability space $(\Omega, \sF, P)$, 
and a Borel measurable function $g: \R^n \to \R^m$ such that $A_g(y) \triangleq \cap_{j=1}^m\set{x \in \R^n: g_j(x) \le y_j} \in \sB(\R^n)$ for all $y \in \R^m$. Then, $g(X): \Omega \to \R^m$ is a random vector. \\
\item <2-> The joint distribution function $F_Y: \R^m \to [0,1]$ for the vector $Y\triangleq g(X)$ is given by 
\EQ{
F_Y(y) = P(X^{-1}(A_g(y))),\quad \text{ for all } y \in \R^m.
} 
\end{itemize}
\end{prop}
\end{frame}

\begin{frame}{Function of random vectors}
\begin{Example}[Sum of random variables]
\begin{itemize}
\item <1-> For a random vector $X: \Omega \to \R^n$ defined on a probability space $(\Omega, \sF, P)$. 
Define an addition function $+: \R^n \to \R$ such that $+(x) = \sum_{i=1}^nx_i$ for any $x \in \R^n$. 
\item <2-> We can verify that $+$ is a Borel measurable function and hence $Y = +(X) = \sum_{i=1}^nX_i$ is a random variable. 
\item <3-> When $n=2$ and $X$ is a continuous random vector with density $f_X: \R^2 \to \R_+$, we can write 
\EQ{
F_Y(y) 
= P(\set{Y \le y}) 
= P(\set{X_1+X_2 \le y}) 
= \int_{x_1 \in \R}\int_{x_2 \le y-x_1}f_X(x_1,x_2)dx_1dx_2.
}
\end{itemize}
\end{Example}
\end{frame}

\begin{frame}{Function of random vectors}
\begin{Example}[Sum of random variables (Continued)]
\begin{itemize}
\item <1-> By applying a change of variable $(x_1, t)= (x_1, x_1 + x_2)$ and changing the order of integration, 
we see that 
\EQ{
F_Y(y) = \int_{t \le y}dt \int_{x_1 \in \R} dx_1 f_X(x_1,t-x_1). 
} 
When $Y$ is a continuous random vector, we can write 
\EQ{
f_Y(y) = \frac{dF_Y(y)}{dy} = \int_{x\in \R}f_X(x, y-x)dx.
}
\end{itemize}
\end{Example}
\end{frame}

\begin{frame}{Function of random vectors}
\begin{Example}[Sum of random variables (Continued)]
\begin{itemize}
\item <1-> When $X: \Omega \to \R^2$ is an independent vector, then $f_X(x) = f_{X_1}(x_1)f_X(x_2)$ for all $x\in\R^2$. 
Therefore, the density of the sum $X_1+X_2$ is given by  
\EQ{
f_Y(y) = \int_{x \in \R}dx f_{X_1}(x)f_{X_2}(y-x) = (f_{X_1}\ast f_{X_2})(y), 
}
where $\ast: \R^{\R} \times \R^{\R} \to \R^{\R}$ is the convolution operator. 
\end{itemize}
\end{Example}
\end{frame}

\begin{frame}{Function of random vectors}
\begin{block}{Theorem}
For a continuous random vector $X:\Omega \to \R^m$ defined on the probability space $(\Omega, \sF, P)$ with density $f_X: \R^m \to \R_+$ and an injective and smooth Borel measurable function $g: \R^m \to \R^m$, 
such that $Y=g(X)$ is a continuous random vector. 
Then the density of random vector $Y$ is given by 
\EQ{
f_Y(y) = \frac{f_X(x)}{\abs{J(y)}},
} 
where $x=g^{-1}(y)$ and $J(y)= (J_{ij}(y) \triangleq \frac{\partial y_j}{\partial x_i}: i,j\in [m])$ is the Jacobian matrix.  
\end{block} 
\end{frame}

\begin{frame}{Function of random vectors}
\begin{Proof}
\begin{itemize}
\item <1-> For an injective map $g: \R^m \to \R^m$ we have $\set{x} = g^{-1}\set{y}$ for any $y \in g(\R^m)$. 
\item <2-> Further, since $g$ is smooth, we have $dy = J(y)dx + o(\abs{dx})$, 
and thus
\EQN{
\label{eqn:Jacobian}
\abs{dy} = \abs{J(y)}\abs{dx} + o(\abs{dx}).
} 
\item <3-> Defining set $dB(y) \triangleq \set{w \in \R^m: y_j \le w_j \le y_j+dy_j}$, 
we observe that for any continuous random vector $Y: \Omega \to \R^m$, we have 
\EQN{
\label{eqn:density}
P \circ Y^{-1} \left(dB(y)\right)= f_Y(y)\abs{dy} = P\circ X^{-1}\Big(dB(x)\Big) 
= f_X(x)\abs{dx}. 
}
\item <3-> We get the result by combining~\eqref{eqn:Jacobian} and~\eqref{eqn:density}.
\end{itemize}
\end{Proof}
\end{frame}


\begin{frame}{Function of random vectors}
\begin{Example}[Sum of random variables]
\begin{itemize}
\item <1-> $X: \Omega \to \R^2$; $Y_1 = X_1 + X_2$. Compute $f_{Y_1}(y_1)$ using the above theorem. 
\item <2-> Define $Y: \Omega \to \R^2$ ; $Y = (X_1+X_2, X_2)$ so that $\abs{J(y)} = 1$. \implies $f_{Y}(y) = f_{X}(x)$.
\item <3-> Thus, we may compute the marginal density of $Y_1$ as,
\EQ{
f_{Y_1}(y_1) 
= \int_{-\infty}^{\infty} f_X(x)\SetIn{x_2= y_2, x_1+x_2 = y_1}dy_2 
= \int_{-\infty}^{\infty} f_X(y_1 -y_2,y_2)dy_2.
}
\item <4-> If $X$ is an independent random vector, then 
\EQ{
f_{Y_1}(y_1) = \int_{-\infty}^{\infty} f_{X_1}(y_1-y_2)f_{X_2}(y_2)dy_2 = (f_{X_1} * f_{X_2})(y_1),
}
where $*$ represents convolution.
\end{itemize}
\end{Example}
\end{frame}
\end{document}
